\section{Method}
\label{sec:method}

\subsection{Tilted-ring model}
\label{sec:ring_model}

We model the stellar disk as a collection of $N=60$ concentric, infinitesimally thin rings spanning radii $R_\mathrm{min}=0.5$\,kpc to $R_\mathrm{max}=18$\,kpc with uniform spacing $\Delta R = (R_\mathrm{max}-R_\mathrm{min})/N \approx 0.29$\,kpc. Each ring $i$ is characterized by its radius $R_i$, mass $M_i$, and angular momentum unit vector $\hat{\boldsymbol{L}}_i(t)$. The direction of $\hat{\boldsymbol{L}}_i$ encodes the instantaneous tilt of ring $i$ relative to the reference plane ($z$-axis).

The ring masses are assigned assuming an exponential surface density profile:
\begin{equation}
    \Sigma(R) = \frac{M_\mathrm{disk}}{2\pi R_d^2} \exp\!\left(-\frac{R}{R_d}\right),
\end{equation}
where $M_\mathrm{disk}=5\times10^{10}\,\mathrm{M}_\odot$ is the total disk mass and $R_d=3.5$\,kpc is the exponential scale length. The mass of ring $i$ is then $M_i = 2\pi R_i \,\Sigma(R_i)\,\Delta R$.

Each ring's circular velocity is computed from the enclosed mass of the combined NFW halo and exponential disk:
\begin{equation}
    v_\mathrm{circ}(R) = \sqrt{\frac{G\left[M_\mathrm{enc}^\mathrm{halo}(R) + M_\mathrm{enc}^\mathrm{disk}(R)\right]}{R}},
\end{equation}
where $M_\mathrm{enc}^\mathrm{disk}(R) = M_\mathrm{disk}\left[1-(1+R/R_d)\exp(-R/R_d)\right]$ and the NFW enclosed mass is given in Section~\ref{sec:halo}. The angular velocity is $\Omega_i = v_\mathrm{circ}(R_i)/R_i$ and the specific angular momentum is $\ell_i = v_\mathrm{circ}(R_i)\,R_i$.

\subsection{Dark matter halo}
\label{sec:halo}

The host galaxy's dark matter halo is modeled as a Navarro--Frenk--White \citep[NFW;][]{navarro1996} profile with virial mass $M_\mathrm{vir}=10^{12}\,\mathrm{M}_\odot$, concentration $c=12$, and virial radius $R_\mathrm{vir}=200$\,kpc. The enclosed mass is
\begin{equation}
    M_\mathrm{enc}^\mathrm{halo}(r) = \frac{M_\mathrm{vir}}{\ln(1+c) - c/(1+c)} \left[\ln(1+x) - \frac{x}{1+x}\right],
\end{equation}
where $x = r/R_s$ and $R_s = R_\mathrm{vir}/c$ is the scale radius.

To model the effect of a slightly oblate halo, we introduce the axis ratio $q \leq 1$, parameterizing the halo flattening. The oblateness produces a torque on tilted rings that causes differential precession. The precession frequency is
\begin{equation}
    \nu_p^2(R) = (1-q^2)\,\frac{G\,M_\mathrm{enc}^\mathrm{halo}(R)}{R^3}.
    \label{eq:nu_p}
\end{equation}
Our fiducial model adopts $q=0.95$, corresponding to a mildly oblate halo.

\subsection{Torques}
\label{sec:torques_method}

The evolution of each ring's angular momentum direction is governed by
\begin{equation}
    \frac{\mathrm{d}\hat{\boldsymbol{L}}_i}{\mathrm{d}t} = \frac{\boldsymbol{T}_i^\mathrm{self} + \boldsymbol{T}_i^\mathrm{halo} + \boldsymbol{T}_i^\mathrm{tidal}}{M_i\,\ell_i},
    \label{eq:dLdt}
\end{equation}
where $M_i\,\ell_i = |\boldsymbol{L}_i|$ is the angular momentum magnitude, and the three torque terms are described below.

\subsubsection{Self-gravity coupling}

When neighboring rings are mutually tilted, they exert gravitational torques on each other that resist differential bending. Following \citet{sparke1988}, the coupling torque on ring $i$ from ring $j$ is
\begin{equation}
    \boldsymbol{T}_{ij}^\mathrm{self} = C_{ij}\,\hat{\boldsymbol{L}}_i \times \left(\hat{\boldsymbol{L}}_j \times \hat{\boldsymbol{L}}_i\right),
\end{equation}
where the coupling coefficient is
\begin{equation}
    C_{ij} = \frac{G\,M_i\,M_j\,\min(R_i,R_j)}{r_\mathrm{eff}^2\,\max(R_i,R_j)},
\end{equation}
and $r_\mathrm{eff} = \sqrt{|R_i-R_j|^2 + h^2}$ is the effective separation softened by the disk scale height $h=z_\mathrm{disk}=0.3$\,kpc. The total self-gravity torque on ring $i$ is $\boldsymbol{T}_i^\mathrm{self}=\sum_{j\neq i} \boldsymbol{T}_{ij}^\mathrm{self}$, summed over the nearest 8 neighbors in each direction.

This coupling provides the disk's \emph{bending stiffness}: it resists differential tilting and enables the propagation of bending waves through the disk.

\subsubsection{Halo torque}

The oblate halo exerts a torque that drives precession of tilted rings around the halo symmetry axis ($\hat{\boldsymbol{z}}$):
\begin{equation}
    \boldsymbol{T}_i^\mathrm{halo} = \nu_p^2(R_i)\,M_i\,\ell_i\,\left(\hat{\boldsymbol{z}} \times \hat{\boldsymbol{L}}_i\right),
    \label{eq:T_halo}
\end{equation}
where $\nu_p^2$ is given by equation~(\ref{eq:nu_p}). This torque is perpendicular to both $\hat{\boldsymbol{z}}$ and $\hat{\boldsymbol{L}}_i$, causing pure precession without changing the tilt angle.

\subsubsection{Tidal torque}
\label{sec:tidal_torque}

The satellite galaxy, treated as a point mass $M_\mathrm{sat}$ at position $\boldsymbol{r}_\mathrm{sat}(t)$, exerts a quadrupole tidal torque on each ring:
\begin{equation}
    \boldsymbol{T}_i^\mathrm{tidal} = \frac{3}{2}\,\frac{G\,M_\mathrm{sat}\,M_i\,R_i^2}{d^3}\,\left(\hat{\boldsymbol{d}} \times \hat{\boldsymbol{L}}_i\right)\left(\hat{\boldsymbol{d}} \cdot \hat{\boldsymbol{L}}_i\right),
    \label{eq:T_tidal}
\end{equation}
where $d = |\boldsymbol{r}_\mathrm{sat}|$ is the distance from the galactic centre to the satellite (softened by $\epsilon_\mathrm{sat}$) and $\hat{\boldsymbol{d}} = \boldsymbol{r}_\mathrm{sat}/d$. The key features of this torque are:
\begin{itemize}
    \item It scales as $R_i^2/d^3$, making the outer disk far more susceptible than the inner disk.
    \item The factor $(\hat{\boldsymbol{d}}\cdot\hat{\boldsymbol{L}}_i)(\hat{\boldsymbol{d}}\times\hat{\boldsymbol{L}}_i)$ vanishes when the satellite is in the disk plane ($\hat{\boldsymbol{d}}\perp\hat{\boldsymbol{L}}_i$) or along the disk normal ($\hat{\boldsymbol{d}}\parallel\hat{\boldsymbol{L}}_i$), and is maximized at an angle of $45^\circ$ between the two vectors.
\end{itemize}

\subsection{Satellite orbit}
\label{sec:orbit}

The satellite follows a fixed Keplerian orbit with semi-major axis $a$, eccentricity $e$, and inclination $i$ relative to the disk plane. The orbital position is obtained by solving Kepler's equation iteratively at each time-step. The pericentric distance is $d_\mathrm{peri} = a(1-e)$ and the orbital period is $P$.

Table~\ref{tab:params} summarizes the fiducial model parameters.

\begin{table}
    \centering
    \caption{Fiducial model parameters.}
    \label{tab:params}
    \begin{tabular}{lcc}
        \toprule
        Parameter & Symbol & Value \\
        \midrule
        \multicolumn{3}{l}{\textit{Stellar disk}} \\
        Total mass & $M_\mathrm{disk}$ & $5\times10^{10}\,\mathrm{M}_\odot$ \\
        Scale length & $R_d$ & 3.5\,kpc \\
        Scale height & $h$ & 0.3\,kpc \\
        Number of rings & $N$ & 60 \\
        Radial range & $R$ & 0.5--18\,kpc \\
        \midrule
        \multicolumn{3}{l}{\textit{NFW dark matter halo}} \\
        Virial mass & $M_\mathrm{vir}$ & $10^{12}\,\mathrm{M}_\odot$ \\
        Concentration & $c$ & 12 \\
        Virial radius & $R_\mathrm{vir}$ & 200\,kpc \\
        Axis ratio & $q$ & 0.95 \\
        \midrule
        \multicolumn{3}{l}{\textit{Satellite galaxy}} \\
        Mass & $M_\mathrm{sat}$ & $10^{11}\,\mathrm{M}_\odot$ \\
        Semi-major axis & $a$ & 25\,kpc \\
        Eccentricity & $e$ & 0.3 \\
        Orbital inclination & $i$ & $45^\circ$ \\
        Orbital period & $P$ & 800\,Myr \\
        Softening length & $\epsilon_\mathrm{sat}$ & 2\,kpc \\
        \midrule
        \multicolumn{3}{l}{\textit{Numerical}} \\
        Time-step & $\Delta t$ & 1\,Myr \\
        Total time & $T$ & 3000\,Myr \\
        Snapshot interval & & 25\,Myr \\
        \bottomrule
    \end{tabular}
\end{table}

\subsection{Numerical integration}
\label{sec:integration}

The system of ordinary differential equations~(\ref{eq:dLdt}) is integrated using a fourth-order Runge--Kutta (RK4) scheme with a fixed time-step $\Delta t = 1$\,Myr. After each full RK4 step, the unit vectors $\hat{\boldsymbol{L}}_i$ are renormalized to maintain $|\hat{\boldsymbol{L}}_i|=1$.

We adopt a unit system in which lengths are in kpc, masses in $10^{10}\,\mathrm{M}_\odot$, and times in Myr. The gravitational constant in these units is $G = 4.498\times10^{-6}\,\mathrm{kpc}^3/(10^{10}\,\mathrm{M}_\odot\,\mathrm{Myr}^2)$.

The total integration time is $T=3000$\,Myr, corresponding to approximately 3.75 orbital periods of the satellite. Snapshots of $\hat{\boldsymbol{L}}_i(t)$ are saved every 25\,Myr for post-processing.

\subsection{Diagnostics}
\label{sec:diagnostics}

We characterize the warp through two primary diagnostics:

\begin{enumerate}
    \item \textbf{Tilt angle} $\theta_i(t) = \arccos(\hat{L}_{i,z})$: the angle between ring $i$'s angular momentum vector and the $z$-axis. A flat disk has $\theta_i=0$ for all $i$.
    
    \item \textbf{Line of nodes} $\phi_i(t) = \arctan(\hat{L}_{i,y}/\hat{L}_{i,x})$: the azimuthal direction of ring $i$'s tilt. A coherent warp has $\phi_i \approx \mathrm{const}$ across all radii; a wound-up warp shows $\phi_i$ varying strongly with $R$.
\end{enumerate}

We also report the maximum tilt angle across all rings, $\theta_\mathrm{max}(t) = \max_i \theta_i(t)$, and the mean tilt of the outer disk, $\langle\theta\rangle_\mathrm{outer}(t) = \langle\theta_i\rangle_{R_i>10\,\mathrm{kpc}}$.
